\chapter{付録}

\section{分析で用いた主な変数の定義（概要）}
本研究で用いた主要変数の概要を、再現性の観点から補足する。

\subsection{リスク開示の新規性（Novelty）}
有価証券報告書「事業等のリスク」欄を一定長のテキスト・チャンクに分割し、当年チャンクと前年チャンク集合の間で埋め込み類似度を計算する。
各当年チャンクについて前年集合に対する最大類似度を求め、\(1-\max(\mathrm{sim})\) をチャンク単位の新規性として定義する。
企業年の新規性指標は、上位 \(k\%\) のチャンク新規性の平均（\texttt{change\_topk}）として集約する。

\subsection{市場反応（アウトカム）}
提出日をイベント日とし、イベント窓（例：\(-1,+1\)）で以下のアウトカムを算出する。

\begin{itemize}
  \item \texttt{abvol\_mkt}：市場調整後の異常出来高（投資家の注目度の代理変数）
  \item \texttt{vol\_abs}：異常収益の絶対値集計（不確実性の代理変数）
  \item \texttt{car\_mm}：市場モデルに基づく累積異常収益（価格反応）
\end{itemize}

\section{再現可能性のための補足（実行手順の最小要約）}
分析はPythonを中心に実施し、データ取得から前処理、指標構築、推定、図表生成までをノートブックで管理した。
また、重い処理（埋め込み計算や疑似イベントの反復推定等）については中間成果物のキャッシュを活用し、乱数シードを固定して再現性を担保した。
実行環境・手順の詳細は、リポジトリ内の \texttt{README} およびノートブックに記載の手順に従う。

\section{追加の頑健性確認（本文の補足）}
本文で示した主結果に加え、以下の観点で追加検証を行った。

\begin{itemize}
  \item 事前トレンド（pre-trend）の確認
  \item 偽イベント日（placebo）による比較
  \item 固定効果（年、業種、企業）の導入有無による結果の安定性
\end{itemize}

これらの追加分析の数値結果や図表は、リポジトリの出力（例：\texttt{outputs/} 配下）として保存しており、主要な含意は本文でも要約した。


